\documentclass[11pt]{article}

\usepackage[margin=1in]{geometry}
\usepackage[T1]{fontenc}
\usepackage[utf8]{inputenc}
\usepackage{lmodern}
\usepackage{setspace}
\usepackage{booktabs}
\usepackage{array}
\usepackage{hyperref}
\usepackage[numbers]{natbib}

\hypersetup{
  colorlinks=true,
  linkcolor=black,
  citecolor=black,
  urlcolor=blue
}

\singlespacing

\title{Compiling Ancient Greek Semantic DAGs from Universal Dependencies}
\author{
Group 7\\
Evan Liu, George Jiang, Bennet Wang, and Haocheng Han\\
New York University
}
\date{May 2026}

\begin{document}
\maketitle

\begin{abstract}
This paper presents a compiler that turns morphosyntactic analyses from Ancient
Greek Universal Dependencies (UD) treebanks into explicit, inspectable, and
structurally valid semantic graphs. Each sentence is converted into a directed
acyclic graph (DAG) whose nodes preserve token, lemma, part of speech, and
morphology, while edges encode a compact semantic inventory:
\textsc{agent}, \textsc{theme}, \textsc{modifier}, \textsc{complement}, and
\textsc{coord}. I evaluate the system in two ways. First, on a manually labeled
historical gold set of 21 real PROIEL and Perseus sentences (236 tokens, 122
gold semantic edges), the complete graph achieves 68.82 precision, 95.90 recall,
and 80.14 F1, compared with 19.55 F1 for a local baseline that ignores UD
dependency edges and predicts semantic relations only from token order, POS, and
case heuristics. Excluding \textsc{modifier} edges raises the score on core
semantic roles to 91.75 precision, 94.68 recall, and 93.19 F1, compared with
22.56 F1 for the same baseline. Second, on the full combined Ancient Greek UD
test material in the repository (2,353 sentences, 31,913 tokens), the
compiler preserves all nodes, produces valid DAGs for every sentence, maps
79.36\% of all non-punctuation dependencies, covers 98.98\% of content
dependencies, and covers 99.96\% of predicate-argument dependencies. The main
limitation is semantic over-generation of modifier edges, especially for
function-like or syntax-only dependents. These results show that UD-backed
graph compilation can recover a strong semantic role structure for
Ancient Greek while preserving transparent graph structure and full DAG validity.
\end{abstract}

\section{Introduction}
Ancient Greek texts are often accessed through translations, commentaries, and
lexical notes, but many research questions require explicit structure over the
original language. A classic example is determining who did what to whom, under
which conditions, and with what modifiers. Dependency treebanks already provide
valuable morphosyntactic information for Ancient Greek, but their labels remain
primarily syntactic. The central question is whether those syntactic analyses can be
compiled into a more semantic, graph-based layer that remains easy to inspect
and evaluate.

The system converts UD-annotated sentences into semantic directed acyclic graphs
(DAGs). Each token becomes one node. Each node preserves the original token
surface form, lemma, part of speech, and morphological features. Syntactic
dependency labels are then mapped into a smaller semantic edge inventory:
\textsc{agent}, \textsc{theme}, \textsc{modifier}, \textsc{complement}, and
\textsc{coord}. The output is structurally validated with graph checks so that
the representation is not only interpretable, but also well formed.

The paper makes four concrete contributions: a compact semantic DAG
representation grounded in real treebank inputs, a compiler from UD annotations
to that graph inventory, a small manual evaluation set of historical Greek
sentences labeled directly in the target format, and an empirical evaluation
that combines baseline comparison with corpus-scale coverage statistics.

\section{Related Work}
This work sits at the intersection of dependency annotation, classical
treebanking, semantic role labeling, and semantic graph representations.

Universal Dependencies provides a cross-linguistically consistent framework for
tokenization, part-of-speech tagging, morphological features, and dependency
relations \citep{nivre2016ud}. In this setting, UD matters because it supplies a
stable annotation layer from which semantic abstraction can begin. This makes it
possible to reorganize syntactic structure into semantic structure while preserving
the original token-level linguistic information.

For Ancient Greek specifically, existing dependency treebanks make this kind of
work possible. The Ancient Greek and Latin Dependency Treebanks described by
\citet{bamman2011agldt} demonstrate the value of large, manually annotated
classical corpora for linguistic and philological research. The PROIEL project
\citep{haug2008proiel} extends the treebanking tradition to parallel and
historical Indo-European texts, including Ancient Greek material that is already
distributed in UD form. These resources provide the morphosyntactic basis for
the system input used here.

At the semantic level, the most relevant precedent is semantic role labeling.
PropBank-style annotation formalizes predicate-argument structure over events and
their participants \citep{palmer2005propbank}. The system adopts the
same general intuition that subject, object, and clausal dependencies can often
be reorganized into semantically meaningful role edges.

The approach is also inspired by graph-based semantic representations such as
Abstract Meaning Representation (AMR), which represents sentence meaning as a
graph over events, entities, and relations \citep{banarescu2013amr}. This
paper takes the same graph-oriented perspective while using a compact inventory
of node types and role labels tailored to Ancient Greek treebank data.

What distinguishes the present system from the work above is its combination of
Ancient Greek UD input, a deliberately small semantic label inventory, explicit
graph validation, and evaluation directly on manually labeled historical Greek
sentences.

\section{Task Definition and Data}
\subsection{Task Definition}
The system input is a sentence with UD-style token-level annotations: surface
tokens, lemmas, part-of-speech tags, morphological features, and dependency
links. The output is a semantic DAG in which every token is preserved as a node
and every retained relation is rewritten into one of five target labels:
\textsc{agent}, \textsc{theme}, \textsc{modifier}, \textsc{complement}, or
\textsc{coord}. The task is to compile these token-level annotations into a
semantic graph.

The graph is intended to support downstream interpretability. A reader should be
able to inspect the output and recover the main event structure, participant
roles, modifier attachments, and clausal relationships without consulting the
original dependency labels. At the same time, the graph should remain
structurally valid: acyclic, internally consistent, and free of malformed edges.

\subsection{Treebank Inputs}
The repository includes two public Ancient Greek UD test files used in the
experiments:

\begin{itemize}
  \item \texttt{data/real\_eval/grc\_perseus-ud-test.conllu}
  \item \texttt{data/real\_eval/grc\_proiel-ud-test.conllu}
\end{itemize}

When converted with the compiler, these files yield 2,353
sentences and 31,913 non-punctuation tokens. These treebanks provide the
morphosyntactic analyses from which the semantic DAGs are constructed.

\subsection{Manual Historical Semantic Gold}
The main semantic evaluation uses
\texttt{data/gold\_semantic/historical\_semantic\_dags.jsonl}. This file contains
21 manually labeled historical Greek sentences selected from the PROIEL and
Perseus UD test materials. The examples total 236 tokens and 122 gold semantic
edges in the target inventory.

This manual gold set is intentionally content-semantic. It emphasizes event
participants, complements, coordination, and semantically relevant modifiers.
Articles, discourse particles, and other purely functional markers are not
systematically rewarded. That choice is important because a naive syntax-to-graph
mapping could otherwise inflate scores by reproducing formal syntactic detail
that does not contribute much to semantic interpretation.

The repository also evaluates a second view of the same manual gold set that
removes \textsc{modifier} from scoring. In that setting, the remaining core roles
(\textsc{agent}, \textsc{theme}, \textsc{complement}, and \textsc{coord}) account
for 94 gold edges. This core-role subset helps isolate whether the system's main
semantic contribution lies in predicate-argument structure or in broad
syntax-to-modifier projection.

\section{Methodology}
\subsection{Node Construction}
The compiler creates one graph node per retained linguistic token. Each node
stores the original surface form, lemma, UD part of speech, morphological
features, a semantic type, a semantic role, and a confidence score. The semantic
typing scheme is intentionally compact. Nouns and proper nouns typically become
\emph{entity} nodes, verbs and auxiliaries become \emph{event} nodes, adjectives
become \emph{property} nodes, adverbs become \emph{manner} nodes, and conjunction
or linker-like items become \emph{connector} nodes. This inventory is designed
to support graph inspection, role attachment, and evaluation across real
historical sentences.

The compiler copies lemma and morphology information directly into the node
representation. This keeps the graph closely aligned with the source treebank
annotations and preserves the linguistic details needed for semantic analysis.

\subsection{Dependency-to-Semantic Mapping}
The core of the system is a rule-based mapping from UD dependency labels into
the target semantic edge inventory. Table~\ref{tab:mapping} summarizes the
main mapping decisions.

\begin{table}[t]
\centering
\caption{Main mapping from UD dependency labels to semantic DAG labels.}
\label{tab:mapping}
\begin{tabular}{>{\raggedright\arraybackslash}p{2.0in} >{\raggedright\arraybackslash}p{1.4in} >{\raggedright\arraybackslash}p{2.5in}}
\toprule
UD labels & DAG label & Intuition \\
\midrule
\texttt{nsubj}, \texttt{csubj}, \texttt{obl:agent} & \textsc{agent} & canonical subjects and agentive obliques \\
\texttt{obj}, \texttt{iobj}, \texttt{nsubj:pass}, \texttt{csubj:pass}, \texttt{obl:arg} & \textsc{theme} & objects and patient-like arguments \\
\texttt{xcomp}, \texttt{ccomp}, \texttt{acl}, \texttt{advcl}, \texttt{advcl:cmp}, \texttt{obl} & \textsc{complement} & clausal and oblique complements \\
\texttt{amod}, \texttt{nmod}, \texttt{nummod}, \texttt{advmod}, \texttt{appos}, \texttt{det} & \textsc{modifier} & attributive, adverbial, or dependent modifiers \\
\texttt{conj} & \textsc{coord} & coordination between linked conjuncts \\
\bottomrule
\end{tabular}
\end{table}

The mapping is intentionally conservative in its label inventory and broad in
coverage. Many syntactic distinctions are collapsed into a small number of
readable semantic edge types. The trade-off is that some syntactically valid
relations, especially modifier-like ones, may become semantically too broad once
projected into the smaller label space.

Several UD relations are intentionally not treated as semantic edges. In the
corpus-scale coverage analysis, the largest unmapped relation types are
\texttt{case}, \texttt{cc}, \texttt{discourse}, \texttt{mark}, and \texttt{cop}.
These are mostly function markers or structural helpers rather than the central
content relations the graph is meant to expose. Excluding them makes the graph
more semantically readable and explains why content-dependency coverage is much
higher than overall dependency coverage.

\subsection{Graph Building and Validation}
After node and edge creation, the compiler builds a directed graph with stable
node identifiers such as \texttt{n0}, \texttt{n1}, and \texttt{n2}. Duplicate
edges, self-loops, and malformed links are removed before validation. Validation
checks confirm that edge labels come from the allowed inventory, endpoints are
valid, and the final graph is acyclic. The implementation also verifies
predicate-argument well-formedness for required structures and uses
\texttt{networkx} for cycle detection.

These checks make the output usable as structured data, not just as a readable
annotation. High local edge accuracy is not enough if the graph as a whole is
invalid.

\section{Experiments}
\subsection{Experiment 1: Manual Historical Gold Evaluation}
The first experiment measures exact edge precision, recall, and F1 on the 21
manually labeled historical sentences. A predicted edge is counted as correct
only if its source token, destination token, and label all match the gold edge.
This metric is strict but appropriate because the output is a graph, not a bag of
independent labels. The experiment also reports DAG validity rate.

Because the system is deterministic and rule-based, no supervised model is
trained and no statistical parameters are tuned. The public UD test files and
manual historical gold set are used only for reported evaluation. Development
work consisted of designing the dependency-to-semantic mapping and validating
the graph construction procedure, rather than fitting rules to a train/dev split.

Two systems are compared on this task. The first is the full compiler described
above, which maps UD dependencies directly into semantic graph edges. The second
is a deliberately simple baseline. The baseline uses the same UD tokens, lemmas,
POS tags, and morphological features, but it ignores UD dependency links. It
instead predicts semantic edges from local word order, POS, and case heuristics
using the repository's existing linear relation predictor. The comparison asks
how much the dependency layer contributes beyond a simple local heuristic.

For both systems, two scores are reported. The \emph{full graph} score evaluates
all five labels. The \emph{core roles} score excludes \textsc{modifier} and
evaluates only \textsc{agent}, \textsc{theme}, \textsc{complement}, and
\textsc{coord}. This split is motivated by early inspection of the outputs: the
compiler appears strong on event structure but less selective about which
syntactic dependents deserve a semantic modifier edge.

\subsection{Experiment 2: Corpus-Scale UD Conversion}
The second experiment measures how the compiler behaves when applied to the full
combined Ancient Greek UD test data in the repository. Here the focus is on
coverage and robustness over the full converted corpus.

The relevant metrics are:
\begin{itemize}
  \item overall dependency coverage: the proportion of non-root,
  non-punctuation UD dependencies that receive a semantic edge;
  \item content dependency coverage: the same calculation, but excluding
  relations such as \texttt{case}, \texttt{cc}, \texttt{mark}, \texttt{cop},
  \texttt{discourse}, \texttt{aux}, \texttt{fixed}, and \texttt{flat:name};
  \item predicate-argument coverage: the proportion of subject, object, and
  argument-like dependencies that are converted into semantic edges;
  \item DAG validity rate: the proportion of converted sentences whose final
  graphs pass structural validation.
\end{itemize}

This experiment asks whether the system scales cleanly to thousands of real
treebank sentences while preserving the main semantic structure.

\subsection{Reproducibility}
The main experimental entry points in the repository are
\texttt{scripts/run\_ud\_dag.py} for corpus-scale conversion and
\texttt{scripts/evaluate\_historical\_gold.py} for manual semantic evaluation.
All quantitative results reported below were regenerated locally from those
scripts inside the cloned project workspace. The project code, data files,
evaluation scripts, and paper sources are available at
\url{https://github.com/Eliguli712/Greek_Learning}.

\section{Results}
Table~\ref{tab:historical-results} shows the main comparison between the full
compiler and the local baseline on the manual historical gold set.

\begin{table}[t]
\centering
\caption{Comparison between the full compiler and the local baseline on 21
Ancient Greek historical sentences.}
\label{tab:historical-results}
\begin{tabular}{lrrrr}
\toprule
System and setting & Precision & Recall & F1 & Gold edges \\
\midrule
Baseline, full graph & 18.06 & 21.31 & 19.55 & 122 \\
Compiler, full graph & 68.82 & 95.90 & 80.14 & 122 \\
Baseline, core roles only & 21.78 & 23.40 & 22.56 & 94 \\
Compiler, core roles only & 91.75 & 94.68 & 93.19 & 94 \\
\bottomrule
\end{tabular}
\end{table}

The comparison is decisive. On the full graph, the compiler improves from 19.55 F1
to 80.14 F1. On core semantic roles, it improves from 22.56 F1 to 93.19 F1.
This gap is large enough that it cannot be explained by minor tuning choices. It
shows that the UD dependency layer is doing real work for the task, especially
for event structure and clausal attachment.

The baseline fails most clearly on complements. Because it ignores dependency
links and relies only on local order and morphology, it never recovers the gold
\textsc{complement} edges in this evaluation. It also performs weakly on
\textsc{agent} and \textsc{theme}, and only modestly on \textsc{coord}. This is
exactly the behavior one would expect from a local heuristic in a language with
richer morphology and less rigid word order than English.

Table~\ref{tab:label-results} breaks down the full compiler's per-label results
on the same historical evaluation set.

\begin{table}[t]
\centering
\caption{Per-label results for the full compiler on the manual historical gold
set.}
\label{tab:label-results}
\begin{tabular}{lrrrr}
\toprule
Label & Precision & Recall & F1 & Gold edges \\
\midrule
AGENT & 80.95 & 100.00 & 89.47 & 17 \\
THEME & 93.33 & 90.32 & 91.80 & 31 \\
MODIFIER & 38.36 & 100.00 & 55.45 & 28 \\
COMPLEMENT & 94.59 & 94.59 & 94.59 & 37 \\
COORD & 100.00 & 100.00 & 100.00 & 9 \\
\bottomrule
\end{tabular}
\end{table}

The strongest result is the compiler's performance on the semantic backbone of
the sentence. Core roles reach 93.19 F1, with particularly strong performance on
\textsc{complement} and perfect performance on \textsc{coord} in this sample.
Both \textsc{agent} and \textsc{theme} also score above 89 F1. These numbers
support the claim that the system is capturing predicate-argument structure
reliably once strong upstream syntax is available.

The weaker result is the full-graph score of 80.14 F1. The gap between full
graph and core roles comes almost entirely from \textsc{modifier}. Modifier
recall is 100\%, which means the compiler almost never misses a gold semantic
modifier once one exists. However, modifier precision is only 38.36\%. In other
words, the system attaches too many dependents as semantic modifiers. This is
consistent with the system design: a broad set of UD syntactic modifiers is
collapsed into one semantic label, so purely structural or lightly semantic
dependents are often over-produced.

Every evaluated historical sentence still yields a valid DAG. The validity rate
is 100\%, which means the structural layer of the system is robust even when the
semantic edge inventory is imperfect.

Table~\ref{tab:coverage-results} presents the corpus-scale conversion results on
the full combined Ancient Greek UD test material.

\begin{table}[t]
\centering
\caption{Corpus-scale conversion statistics on the combined Ancient Greek UD
test files in the repository.}
\label{tab:coverage-results}
\begin{tabular}{lrr}
\toprule
Metric & Value & Notes \\
\midrule
Sentences converted & 2,353 & Perseus + PROIEL test material \\
Tokens converted & 31,913 & non-punctuation tokens \\
Mapped dependencies & 23,458 & semantic edges produced \\
Unmapped dependencies & 6,102 & mostly function markers \\
Overall dependency coverage & 79.36\% & all non-root, non-punctuation deps \\
Content dependency coverage & 98.98\% & excluding function-marker relations \\
Predicate-argument coverage & 99.96\% & subject/object/argument-style deps \\
DAG validity rate & 100.00\% & all converted graphs valid \\
\bottomrule
\end{tabular}
\end{table}

These coverage results sharpen the interpretation of the manual gold numbers.
The system's lower full-graph precision is not caused by a failure to preserve
predicate-argument structure. In fact, predicate-argument coverage is essentially
complete at 99.96\%. Instead, the remaining gap comes from exactly the kinds of
dependencies the compiler treats cautiously: function markers and broad
modifier-like relations.

The unmapped dependency counts support that interpretation. The largest unmapped
categories are \texttt{case} (2,322), \texttt{cc} (1,660), \texttt{discourse}
(788), \texttt{mark} (580), and \texttt{cop} (490). These are mostly structural
or functional relations rather than the main semantic relations the system is
trying to expose. Once such categories are excluded, content dependency coverage
rises to 98.98\%.

\section{Discussion}
At its core, the paper addresses a clear technical problem: compiling an existing
morphosyntactic analysis into a semantic graph for Ancient Greek. The evaluation
shows that this formulation supports strong predicate-argument recovery,
transparent graph structure, and reliable large-scale conversion across real
treebank data.

The strongest outcome is that a compact semantic label inventory is enough to
recover the predicate-argument backbone of real historical Greek sentences. The
manual evaluation shows this directly, and the corpus-scale coverage analysis
supports it from another angle. The graph representation is therefore a useful
way to expose role structure, clausal structure, and modifier attachment in a
form that remains easy to inspect and validate.

The main weakness is modifier over-generation. Several design decisions
contribute to this. First, the mapping groups many syntactically different UD
relations under one semantic label. Second, the historical gold deliberately
omits many function-like or low-content modifiers, so a syntactically faithful
projection can look semantically noisy. Third, Ancient Greek often contains
particles, determiners, and other dependents whose semantic contribution is
subtle or context dependent. A single \textsc{modifier} label is too coarse to
handle all of these cases well.

The main limitations extend beyond modifiers. The manual semantic
gold set is small, and the system inherits the quality of the upstream UD
annotations used as input. Even so, the results are strong and consistent. The
baseline comparison shows that the final compiler captures much more of the
semantic structure than a local heuristic, and the main remaining error source
is precise semantic filtering of modifiers rather than event structure itself.
These results should not be read as a direct state-of-the-art comparison with
full Ancient Greek parsers, semantic role labelers, or AMR-style systems. The
project instead contributes a reproducible compiler and evaluation setup that
could later be paired with stronger upstream parsing or learned semantic
filtering.

\section{Conclusion and Future Work}
The paper presents a semantic DAG compiler for Ancient Greek. The system preserves token,
lemma, part of speech, morphology, and syntactic structure, then reorganizes
those inputs into a compact semantic graph with explicit role labels. On 21
manually labeled historical sentences, the compiler reaches 80.14 F1 on the full
graph and 93.19 F1 on core semantic roles, compared with 19.55 and 22.56 F1 for
a local baseline that ignores dependency edges. On the full combined UD test
files in the repository, it converts 2,353 sentences with 100\% DAG validity and
near complete coverage of predicate-argument dependencies.

Future work should focus on the part of the system that the evaluation exposed as
fragile: modifier filtering. One direction is to split \textsc{modifier} into a
finer taxonomy that distinguishes semantic attributes, determiners, discourse
markers, and syntax-only attachments. Another is to learn a pruning or
relabeling layer from manually annotated semantic graphs instead of relying on a
single deterministic projection. A third is to expand the manual gold set so
that evaluation covers more authors, genres, and clause types. Longer-term work
could also integrate richer semantic information, such as valency classes,
animacy, or limited coreference, while preserving the graph validity guarantees
that made the current system attractive.

The result is a reproducible compiler that turns real Ancient Greek treebank
analyses into inspectable graph structure.

\section*{Contributions}
Evan Liu led the final evaluation and reporting work, including the UD-backed DAG
evaluation, historical semantic gold evaluation, baseline comparison,
presentation materials, and final paper preparation; he also delivered the final
presentation. Haocheng Han contributed the initial project setup, morpheme and
lexeme framework, lexicon resources, and early design documentation. Bennet Wang
contributed the minimal viable compiler pipeline, phoneme processing,
validation/evaluation utilities, and Aristotle test outputs. George Jiang
contributed updates to the semantic tokenization component.

\bibliographystyle{plainnat}
\bibliography{references}

\end{document}
